\chapter{Síntese: uma arquitetura de referência}
\label{cap:sintese}

\begin{objetivos}
Ao final deste capítulo, você será capaz de:
\begin{itemize}[leftmargin=*,itemsep=0.2em]
  \item reconstruir a trajetória metodológica completa e identificar o que produziu cada
        salto de desempenho;
  \item consolidar resultados heterogêneos em uma visão comparável;
  \item derivar, a partir da evidência, qual componente deve ser usado para qual tarefa;
  \item avaliar criticamente uma arquitetura de referência que integre modelos numéricos
        e camada de linguagem;
  \item enunciar os princípios de projeto que o percurso deste livro sustenta.
\end{itemize}
\end{objetivos}

\section{Onde chegamos}

Este livro contou uma história em quatro atos.

No primeiro, um sistema dual --- regras e autoencoder --- entrou em operação em mais de
duzentos e cinquenta poços e detectou eventos que ninguém havia catalogado. Ele
funcionava, e não sabia dizer o que via.

No segundo, o ensemble de classificadores binários trocou a pergunta ``qual das nove
classes?'' pela pergunta ``é esta? sim ou não'', e com isso tornou expressável a resposta
``nenhuma delas''. O experimento da classe escondida deu $89\,\%$.

No terceiro, duas mudanças conceituais --- segmentar em vez de detectar, e representar
com um classificador em vez de um autoencoder --- levaram a detecção de novidade de
$52\,\%$ a $98{,}2\,\%$.

No quarto, uma camada de linguagem recebeu as estatísticas do segmento e produziu
hipóteses justificadas, reconstruindo descrições fisicamente corretas de eventos cujo
nome nunca lhe foi apresentado.

Este capítulo consolida o percurso e propõe a arquitetura que ele sustenta.

\section{O que produziu cada ganho}

\begin{table}[htbp]
\centering
\caption{Marcos de desempenho e a mudança conceitual responsável por cada um.}
\label{tab:sint-marcos}
\small
\begin{tabular}{p{2.1cm}p{4.3cm}p{2.4cm}p{3.2cm}}
\toprule
\textbf{Capítulo} & \textbf{Mudança conceitual} & \textbf{Resultado} &
\textbf{Custo} \\
\midrule
\cref{cap:dual} &
Combinar regras determinísticas com aprendizado não supervisionado &
ACC $0{,}9727 \to 0{,}9894$ &
Duas bases de código \\
\addlinespace
\cref{cap:binarios} &
Escores independentes em vez de \ing{softmax} acoplada &
Rejeição torna-se expressável &
$9\times$ o treinamento \\
\addlinespace
\cref{cap:closedworld} &
Validar a rejeição com classe escondida &
$89\,\%$ de rejeição &
Nenhum \\
\addlinespace
\cref{cap:mundoaberto} &
Organizar o rejeitado em classes virtuais &
$80{,}8\,\%$ de agrupamento &
Muita engenharia \\
\addlinespace
\cref{cap:segmentacao} &
Segmentar em vez de detectar &
ACC $0{,}500 \to 0{,}863$; desvio $0{,}187 \to 0{,}082$ &
Anotação por amostra \\
\addlinespace
\cref{cap:mahalanobis} &
Representar com classificador, não com autoencoder &
Novidade $52 \to 98{,}2\,\%$ &
Trocar uma linha \\
\addlinespace
\cref{cap:resultados-agente} &
Descrever em linguagem natural &
$89{,}7\,\%$ de novidade; nomes fisicamente corretos &
Serviço externo \\
\bottomrule
\end{tabular}
\end{table}

\begin{destaque}[Os dois maiores ganhos foram os mais baratos]
Trocar o autoencoder pelo latente do classificador multiclasse elevou a detecção de
novidade em quarenta e seis pontos percentuais. Em termos de implementação, é extrair a
saída de uma camada diferente.

Segmentar em vez de detectar elevou a acurácia em trinta e seis pontos e reduziu a
variância entre poços em mais da metade. Exigiu uma arquitetura nova, mas nenhuma
sofisticação --- a U-Net é de 2015.

Em contraste, o ciclo de mundo aberto do \cref{cap:mundoaberto} --- estimação automática
de $K$, interseção entre dois algoritmos de agrupamento, limiar adaptativo com três
regimes, salvaguarda de dois terços --- exigiu esforço de engenharia desproporcional ao
seu resultado, porque operava sobre um espaço sem estrutura de classe.

A lição é geral e vale ser lembrada: \textbf{antes de sofisticar o método, verifique o
espaço em que ele opera}.
\end{destaque}

\section{Resultados consolidados}

\subsection{Detecção de novidade}

\begin{table}[htbp]
\centering
\caption{Detecção de novidade ao longo do livro. Valores em porcentagem. Traços indicam
que a combinação não foi avaliada.}
\label{tab:sint-novidade}
\small
\begin{tabular}{lcccc}
\toprule
& \textbf{Ens.\ bin.} & \textbf{Mahal.} & \textbf{Mahal.} & \textbf{Agente} \\
\textbf{Classe} & \textbf{latente AE} & \textbf{bruto} & \textbf{latente MC} & \textbf{LLM} \\
\midrule
1 & $22{,}8$ & $41{,}5$ & $99{,}0$ & $70{,}8$ \\
2 & $76{,}9$ & $77{,}1$ & $96{,}5$ & $100{,}0$ \\
3 & $48{,}1$ & $61{,}0$ & $99{,}0$ & $83{,}7$ \\
4 & $44{,}7$ & $57{,}8$ & $96{,}0$ & --- \\
5 & $0{,}1$  & $62{,}3$ & $67{,}5$ & --- \\
6 & $47{,}6$ & $47{,}6$ & $99{,}0$ & $95{,}7$ \\
7 & $36{,}5$ & $36{,}7$ & $52{,}3$ & $94{,}9$ \\
8 & $50{,}5$ & $91{,}1$ & $99{,}0$ & $91{,}8$ \\
9 & $51{,}9$ & $88{,}7$ & $99{,}0$ & $54{,}5$ \\
ICV & $52{,}0$ & $67{,}5$ & $98{,}2$ & --- \\
\midrule
\textbf{Global} & --- & --- & $\bm{95{,}4}$ & $89{,}7$ \\
\bottomrule
\end{tabular}
\end{table}

\begin{destaque}[A linha mais interessante da tabela é a da classe 7]
A incrustação no PCK é o pior caso do Mahalanobis latente ($52{,}3\,\%$) --- e um dos
melhores do agente ($94{,}9\,\%$).

A explicação está na natureza da evidência. O método geométrico enxerga uma janela de
duzentos passos e não encontra desvio suficiente, porque a incrustação se desenvolve ao
longo de semanas. O agente recebe a \emph{tendência} --- a inclinação, a razão de
oscilação, a comparação com o regime de base --- que já é uma síntese temporal, e
reconhece o padrão.

Não é que um método seja melhor. É que eles falham em situações diferentes, e isso é
precisamente o que torna a combinação valiosa.
\end{destaque}

Vale registrar a simetria: para a classe 9, a relação se inverte --- $99{,}0\,\%$ do
Mahalanobis contra $54{,}5\,\%$ do agente. A complementaridade é bidirecional.

\subsection{Classificação}

\begin{table}[htbp]
\centering
\caption{Acurácia de classificação por abordagem.}
\label{tab:sint-classificacao}
\small
\begin{tabular}{llc}
\toprule
\textbf{Capítulo} & \textbf{Abordagem} & \textbf{Acurácia global} \\
\midrule
\cref{cap:binarios} & Ensemble binário (espaço bruto)      & $0{,}87$ \\
\cref{cap:binarios} & Multiclasse (espaço bruto)           & $0{,}86$ \\
\cref{cap:mahalanobis} & Binários no latente do multiclasse & $0{,}755 \pm 0{,}171$ \\
\cref{cap:mundoaberto} & Binários no latente do autoencoder & $0{,}730$ \\
\cref{cap:binarios} & OCSVM (espaço bruto)                 & $0{,}65$ \\
\cref{cap:resultados-agente} & Agente, top-3               & $0{,}639$ \\
\cref{cap:resultados-agente} & Agente, top-1               & $0{,}351$ \\
\bottomrule
\end{tabular}
\end{table}

A ordenação é inequívoca: para classificar o conhecido, os métodos supervisionados
clássicos dominam. O agente fica em último lugar, e a distância não é pequena.

\subsection{As duas assimetrias}

Duas comparações resumem o livro inteiro.

\begin{center}
\small
\begin{tabular}{lcc}
\toprule
& \textbf{Afirmar identidade} & \textbf{Reconhecer incompatibilidade} \\
\midrule
Métodos numéricos & $87\,\%$   & $95{,}4\,\%$ \\
Agente de linguagem & $35{,}1\,\%$ & $89{,}7\,\%$ \\
\bottomrule
\end{tabular}
\end{center}

\begin{destaque}[Duas leituras de uma mesma tabela]
\textbf{Por coluna}: reconhecer incompatibilidade é mais fácil que afirmar identidade,
para \emph{ambas} as famílias de método. Isso é uma propriedade do problema. Basta uma
evidência contraditória para descartar uma hipótese; afirmar identidade exige
convergência de todas as evidências.

\textbf{Por linha}: os métodos numéricos são melhores nas duas tarefas, mas a diferença é
de cinquenta e dois pontos percentuais na primeira e de apenas seis na segunda. É onde a
distância é pequena que a contribuição do agente pode ser aproveitada --- somada ao fato
de que ele, e apenas ele, produz uma descrição.
\end{destaque}

\section{Divisão de trabalho por classe}

A \cref{tab:sint-papeis} atribui a cada classe o papel em que cada componente é
efetivamente bom, com base na evidência acumulada.

\begin{table}[htbp]
\centering
\caption{Divisão de trabalho derivada dos resultados. ``Numérico'' refere-se ao pipeline
U-Net + latente do multiclasse + Mahalanobis.}
\label{tab:sint-papeis}
\small
\begin{tabular}{clp{4.4cm}p{4.0cm}}
\toprule
\textbf{Cl.} & \textbf{Evento} & \textbf{Papel do numérico} & \textbf{Papel do agente} \\
\midrule
1 & BSW & Decisor ($99{,}0\,\%$) & Segunda opinião; nomear novidade \\
\addlinespace
2 & DHSV & Decisor ($96{,}5\,\%$) & Confirmar, justificar e validar ($95\,\%$) \\
\addlinespace
3 & Golfadas & Decisor ($99{,}0\,\%$) & Segunda opinião; nomear novidade \\
\addlinespace
4 & Instabilidade & \textbf{Sinalizador}, não classificador & Reordenar as três candidatas para o operador \\
\addlinespace
5 & Perda de produtiv. & \textbf{Sinalizador}, não classificador & Reordenar as três candidatas para o operador \\
\addlinespace
6 & Restrição PCK & Decisor ($99{,}0\,\%$) & Segunda opinião; nomear novidade \\
\addlinespace
7 & Incrustação & Fraco ($52{,}3\,\%$) & \textbf{Decisor auxiliar} ($94{,}9\,\%$); rejeita $100\,\%$ dos erros \\
\addlinespace
8 & Hidrato prod. & Decisor ($99{,}0\,\%$) & Segunda opinião; nomear novidade \\
\addlinespace
9 & Hidrato serv. & Decisor ($99{,}0\,\%$) & \textbf{Cético}: acionar apenas para rejeitar ($88\,\%$) \\
\bottomrule
\end{tabular}
\end{table}

\begin{destaque}[Esta tabela não foi decidida; foi medida]
Nenhuma das atribuições acima foi postulada a priori. Cada uma decorre de um número
reportado nos capítulos anteriores.

Isso é diferente da prática comum de propor uma arquitetura e depois avaliá-la. Aqui a
avaliação por classe --- que exigiu resistir à tentação de reportar apenas médias
globais --- produziu a arquitetura.

É o argumento mais forte deste livro a favor de reportar desempenho desagregado.
\end{destaque}

\subsection{As classes sintomáticas}

As classes 4 e 5 recebem, na \cref{tab:sint-papeis}, um tratamento distinto de todas as
demais: elas deixam de ser classes.

A recomendação é que o sistema, ao identificar um segmento compatível com elas, emita
\emph{``há instabilidade de fluxo; investigue a causa''} em vez de \emph{``esta é a classe
4''}. A primeira saída é verdadeira e útil; a segunda é uma afirmação sobre uma categoria
que não é internamente coerente.

O suporte empírico é a consistência com que essas classes falham: $0{,}69$ e $0{,}79$ nos
ensembles binários; $67{,}5\,\%$ de detecção de novidade para a classe 5 contra mais de
$96\,\%$ das demais; $18{,}5\,\%$ e $26{,}7\,\%$ de acerto do agente. Quatro métodos
independentes, o mesmo resultado.

\begin{destaque}[Uma recomendação para o próprio conjunto de dados]
A conclusão que este livro sustenta é que as classes 4 e 5 deveriam ser
\emph{redefinidas} no 3W --- por mecanismo, e não por sintoma --- ou explicitamente
marcadas como categorias sintomáticas, para as quais métricas de classificação não são
apropriadas.
\end{destaque}

\section{A arquitetura de referência}

\subsection{Estrutura em camadas}

\begin{enumerate}[leftmargin=*,itemsep=0.4em]

\item \textbf{Aquisição e validação.} Diagrama de decisão do \cref{cap:dual}:
identificação do poço, validação do dado, discriminação de transiente. Determinístico,
auditável, funciona no primeiro segundo.

\item \textbf{Segmentação.} U-Net unidimensional por sensor (\cref{cap:segmentacao}).
Produz a delimitação de que todas as camadas seguintes dependem.

\item \textbf{Representação.} Penúltima camada do classificador multiclasse
(\cref{cap:mahalanobis}). Sessenta e quatro dimensões com estrutura de classe.

\item \textbf{Decisão conhecido/novidade.} Distância de Mahalanobis com limiar no
percentil 95.

\item \textbf{Classificação.} Ensemble binário sobre o latente, para os segmentos
julgados conhecidos.

\item \textbf{Interpretação.} Métricas do \cref{cap:agente} e agente de linguagem, para
todos os segmentos anômalos.

\item \textbf{Decisão.} Operador humano, com as saídas das camadas 4 a 6 apresentadas
separadamente.

\item \textbf{Incorporação.} Ciclo de classes virtuais (\cref{cap:mundoaberto}), agora
sobre um espaço latente com estrutura, e com nomes propostos pelo agente.

\end{enumerate}

\subsection{Fluxo}

\begin{center}
\small
\begin{tabular}{ll}
\toprule
\textbf{Situação} & \textbf{Comportamento do sistema} \\
\midrule
Dado inválido & Descartado na camada 1; não gera alerta \\
Transiente conhecido & Descartado na camada 1 \\
Normal & Camada 4 não sinaliza; registro apenas \\
Anomalia conhecida & Camadas 5 e 6; agente valida ou contesta \\
Anomalia desconhecida & Camada 4 sinaliza; agente nomeia; camada 8 acumula \\
Classes 4 ou 5 & Sinalizador; três candidatas do agente ao operador \\
\bottomrule
\end{tabular}
\end{center}

\begin{destaque}[A camada 8 fecha o ciclo, agora com nome]
O \cref{cap:mundoaberto} terminou com classes virtuais anônimas --- ``grupo 11'' --- e
apontou isso como sua principal limitação.

Nesta arquitetura, a camada 8 recebe segmentos rejeitados por um método que acerta
$95{,}4\,\%$, agrupados em um espaço com silhueta positiva, e acompanhados de descrições
geradas pelo agente com taxa de convergência mensurável. O especialista recebe: um grupo
coeso, um nome proposto, uma justificativa ancorada em estatísticas verificáveis, e um
indicador de quão consistente foi a descrição.

Confirmar ou corrigir isso é trabalho de minutos. Era trabalho de horas.
\end{destaque}

\section{Sete princípios}

O percurso deste livro sustenta sete afirmações que transcendem o problema de poços
submarinos.

\begin{enumerate}[leftmargin=*,itemsep=0.5em]

\item \textbf{A representação determina o que é detectável.} Um mesmo escore de
distância detecta $67{,}5\,\%$ ou $98{,}2\,\%$ da mesma anomalia, conforme o espaço.
Antes de trocar o método, verifique o espaço.

\item \textbf{Reconhecer incompatibilidade é mais fácil que afirmar identidade.} Vale
para métodos geométricos e para modelos de linguagem. Projete sistemas que explorem essa
assimetria em vez de ignorá-la.

\item \textbf{Escores acoplados não podem expressar rejeição.} A restrição
$\sum_k \hat{p}_k = 1$ é uma decisão de projeto com consequência epistemológica, e quase
sempre é tomada sem que se note.

\item \textbf{Reportar desempenho desagregado é o que permite projetar arquiteturas.} A
divisão de trabalho da \cref{tab:sint-papeis} é invisível em qualquer média global.

\item \textbf{Reduzir a variância vale mais que aumentar a média.} Um sistema com
acurácia $0{,}86 \pm 0{,}08$ é implantável; um com $0{,}86 \pm 0{,}19$ não é, porque o
operador não sabe quando confiar.

\item \textbf{Resultados negativos são informação.} O empate entre autoencoder, VAE e DCN
no \cref{cap:mundoaberto} foi o que apontou o gargalo que o \cref{cap:mahalanobis}
resolveu.

\item \textbf{Um sistema que descreve o que vê audita o conjunto de dados.} O artefato de
telemetria do 3W só se tornou visível quando quatro classes receberam nomes contendo
``perda de telemetria''. Nenhuma matriz de confusão o revelou.

\end{enumerate}

\resumocapitulo{%
A trajetória do livro produziu ganhos cuja magnitude foi inversamente proporcional ao seu
custo de implementação: trocar o autoencoder pelo latente do classificador multiclasse
elevou a detecção de novidade em quarenta e seis pontos percentuais e custa extrair a
saída de outra camada, enquanto o elaborado ciclo de agrupamento do
\cref{cap:mundoaberto} exigiu esforço desproporcional por operar sobre um espaço sem
estrutura de classe. Consolidados, os resultados revelam duas assimetrias: reconhecer
incompatibilidade é mais fácil que afirmar identidade, para ambas as famílias de método;
e a vantagem dos métodos numéricos sobre o agente é de cinquenta e dois pontos na
classificação e de apenas seis na detecção de novidade. A avaliação desagregada por
classe permitiu derivar --- e não postular --- uma divisão de trabalho: métodos numéricos
decidem nas classes 1, 2, 3, 6, 8 e 9; o agente é decisor auxiliar na classe 7, onde o
método geométrico atinge apenas $52{,}3\,\%$; é cético na classe 9; e as classes
sintomáticas 4 e 5 devem ser tratadas como sinalizadores, não como classes. A arquitetura
de referência resultante tem oito camadas, da validação determinística do dado ao ciclo
de incorporação de classes virtuais --- que agora recebe grupos coesos, nomes propostos e
justificativas verificáveis. Sete princípios generalizáveis encerram o capítulo.}

\begin{exercicios}
\item A \cref{tab:sint-novidade} mostra que a classe 7 é o pior caso do método geométrico
      e um dos melhores do agente, enquanto a classe 9 apresenta a relação inversa.
      Formule uma regra de seleção de método por classe e estime o desempenho global de
      um sistema que a aplique.

\item Implemente a arquitetura de referência de oito camadas e meça o desempenho
      fim-a-fim sobre o 3W. Compare com cada componente isolado. O ganho da integração é
      aditivo, subaditivo ou superaditivo?

\item O princípio 5 afirma que reduzir variância vale mais que aumentar média. Formalize
      essa intuição: sob que função de utilidade do operador ela é verdadeira? Construa um
      contraexemplo em que ela seja falsa.

\item Proponha uma redefinição das classes 4 e 5 do 3W baseada em mecanismo em vez de
      sintoma. Que informação adicional seria necessária para rotular as instâncias
      existentes sob a nova definição?

\item A arquitetura é sequencial e erros se propagam. Estime a acurácia fim-a-fim em
      função das acurácias das camadas 2, 4 e 6, sob independência. Compare com o valor
      medido no exercício 2 e discuta a diferença.

\item Projete um mecanismo de realimentação em que as correções do operador sobre as
      classes virtuais melhorem as camadas 3 e 5. Que riscos de deriva esse laço
      introduz?

\item Escolha um dos sete princípios e teste-o em um domínio completamente diferente ---
      diagnóstico médico, detecção de fraude, manutenção de equipamentos rotativos.
      Ele se sustenta? Que adaptação exige?
\end{exercicios}

\begin{leituras}
\item \citet{LopesTese2025}, \citet{Lopes2026PST}, \citet{Lopes2026OTC} e
      \citet{Lopes2026LLM} --- os quatro trabalhos consolidados neste capítulo.
\item \citet{Scheirer2013} --- o arcabouço teórico de mundo aberto que atravessa todo o
      percurso.
\item \citet{VazVargas2026_3w2} --- a evolução do 3W, relevante para a recomendação sobre as
      classes sintomáticas.
\end{leituras}
